\documentclass[UTF8,zihao=-4,fontset=fandol]{ctexart}
\usepackage[a4paper,margin=2.15cm,headheight=15pt]{geometry}
\usepackage{fontspec}
\usepackage{xeCJK}
\usepackage{xcolor}
\usepackage{hyperref}
\usepackage{bookmark}
\usepackage{titlesec}
\usepackage{fancyhdr}
\usepackage{array}
\usepackage{longtable}
\usepackage{colortbl}
\usepackage[most]{tcolorbox}
\usepackage{enumitem}
\usepackage{microtype}

\definecolor{ink}{HTML}{111827}
\definecolor{muted}{HTML}{64748B}
\definecolor{accent}{HTML}{1F3A5F}
\definecolor{accentLight}{HTML}{EEF3F8}
\definecolor{headerBg}{HTML}{F1F5F9}
\definecolor{line}{HTML}{CBD5E1}
\definecolor{soft}{HTML}{F8FAFC}

\hypersetup{
  colorlinks=true,
  linkcolor=accent,
  urlcolor=accent,
  citecolor=accent,
  pdftitle={智汇大学英语听说教程4 Theme 2 中英对照文字稿},
  pdfauthor={Codex},
  pdfsubject={WE Learn transcript bilingual notes}
}

\pagestyle{fancy}
\fancyhf{}
\lhead{\textcolor{muted}{智汇大学英语听说教程4 / Theme 2}}
\rhead{\textcolor{muted}{中英对照文字稿}}
\cfoot{\thepage}
\renewcommand{\headrulewidth}{0.4pt}
\arrayrulecolor{line}

\titleformat{\section}
  {\Large\bfseries\color{accent}}
  {\thesection}
  {0.75em}
  {}
\titleformat{\subsection}
  {\large\bfseries\color{ink}}
  {\thesubsection}
  {0.75em}
  {}
\titlespacing*{\section}{0pt}{1.2em}{0.6em}

\setlength{\parindent}{0pt}
\setlength{\parskip}{0.72em}
\setlist{nosep,leftmargin=1.4em}
\hyphenpenalty=10000
\exhyphenpenalty=10000
\emergencystretch=2em
\newcommand{\TranscriptColumn}{\raggedright\arraybackslash\hyphenpenalty=10000\exhyphenpenalty=10000\sloppy}

\newtcolorbox{infobox}{
  enhanced,
  breakable,
  colback=accentLight,
  colframe=accent,
  boxrule=0.45pt,
  arc=2mm,
  left=8pt,
  right=8pt,
  top=6pt,
  bottom=6pt,
  before upper={\raggedright}
}

\begin{document}
\begin{titlepage}
\centering
\vspace*{2.9cm}
{\Huge\bfseries\textcolor{accent}{智汇大学英语听说教程4}\par}
\vspace{0.45cm}
{\LARGE Theme 2 中英对照视频文字稿\par}
\vspace{0.55cm}
{\textcolor{line}{\rule{0.52\textwidth}{0.6pt}}\par}
\vfill
{\large\textcolor{muted}{生成日期：2026-06-02}\par}
\end{titlepage}

\pagenumbering{roman}
\tableofcontents
\clearpage
\pagenumbering{arabic}


\clearpage
\section{Session 2.1 Activity 4 \& 5: AGI, the Top One Emerging Technology that Are Changing Our World}
\addcontentsline{toc}{subsection}{中文译名：改变世界的顶尖新兴技术：人工通用智能}

\renewcommand{\arraystretch}{1.23}
\begin{longtable}{>{\TranscriptColumn}p{0.47\textwidth} >{\TranscriptColumn}p{0.47\textwidth}}
\rowcolor{headerBg}\textbf{English Transcript} & \textbf{中文翻译} \\ \hline
\textbf{01}\quad Technologies are changing our world forever, but not in ways you might expect. For example, artificial general intelligence. Artificial general intelligence, often referred to as AGI, is the concept of a machine with the ability to understand, learn, and apply its intelligence to solve any problem, much like a human being. Unlike narrow AI, which is designed for specific tasks, AGI would have the capacity to handle a wide range of cognitive tasks and adapt to new situations autonomously. & \textbf{01}\quad 技术正在永久地改变我们的世界，但改变的方式可能不是你预想的那样。比如人工通用智能。人工通用智能常被称为 AGI，指的是机器具有像人一样理解、学习并运用智能来解决任何问题的能力。不同于只为特定任务设计的窄域 AI，AGI 将有能力处理范围很广的认知任务，并能自主适应新的情境。 \\ \hline
\textbf{02}\quad Significant strides are being made in AGI research by leading organizations such as OpenAI and Google DeepMind. One approach is through deep learning and neural networks, which mimic the human brain structure and function. Researchers are working on expanding these models to handle more complex abstract tasks that go beyond pattern recognition and data processing. Another interesting development is in the field of reinforcement learning, where AI systems learn to make decisions by trial and error, receiving rewards for successful outcomes. This approach is seen as a potential pathway towards developing more general problem-solving capabilities in AI. & \textbf{02}\quad OpenAI、Google DeepMind 等领先机构正在 AGI 研究方面取得显著进展。一种路径是深度学习和神经网络，它们模仿人脑的结构和功能。研究人员正在扩展这些模型，使它们能够处理超出模式识别和数据处理的、更复杂、更抽象的任务。另一个有意思的发展在强化学习领域：AI 系统通过试错学习如何做决定，并在成功结果之后获得奖励。这种方法被看作是发展 AI 更通用的问题解决能力的一条潜在路径。 \\ \hline
\textbf{03}\quad Looking into the future, the evolution of AGI could have profound implications. If achieved, AGI could perform a wide range of tasks, from complex scientific research and medical diagnosis to creative arts and decision-making. We may arrive at a stage where AGI is capable of handling any task that involves computer-related work, and it could potentially surpass human intelligence in certain areas, leading to breakthroughs in various fields. & \textbf{03}\quad 展望未来，AGI 的演进可能产生深远影响。如果 AGI 得以实现，它可以执行大量任务，从复杂的科学研究、医学诊断，到创造性艺术和决策制定。我们可能会进入这样一个阶段：AGI 能够处理任何涉及计算机相关工作的任务，并且在某些领域可能超过人类智能，从而带来各个领域的突破。 \\ \hline
\textbf{04}\quad AGI could also lead to the development of more intuitive and versatile personal assistance, capable of understanding and responding to a wide range of human needs and preferences. In industry, AGI could automate complex tasks, leading to significant efficiency gains. In time, artificial general intelligences may have the capability to enhance their own algorithms and architectures, potentially giving rise to super-intelligent AIs. These super-intelligent AIs could possess intelligence that surpasses human capabilities by thousands or even millions of times. Such advanced AIs hold the potential to create groundbreaking technologies and change society in ways that are currently difficult to fully comprehend. & \textbf{04}\quad AGI 还可能促成更直观、更多用途的个人助手，它们能够理解并回应人类多种多样的需求和偏好。在工业领域，AGI 可以自动化复杂任务，带来显著的效率提升。随着时间推移，人工通用智能可能具备增强自身算法和架构的能力，从而可能产生超级智能 AI。这些超级智能 AI 的智能可能超过人类能力数千倍甚至数百万倍。这样的高级 AI 有潜力创造突破性技术，并以我们目前难以完全理解的方式改变社会。 \\ \hline
\end{longtable}


\clearpage
\section{Session 2.1 Activity 7: How Blockchain Works (Part 1)}
\addcontentsline{toc}{subsection}{中文译名：区块链如何运作（上）}

\renewcommand{\arraystretch}{1.23}
\begin{longtable}{>{\TranscriptColumn}p{0.47\textwidth} >{\TranscriptColumn}p{0.47\textwidth}}
\rowcolor{headerBg}\textbf{English Transcript} & \textbf{中文翻译} \\ \hline
\textbf{01}\quad It's called the blockchain. Blockchain. Now it's not the most sonorous word in the world, but I believe that this is now the next generation of the internet, and that it holds vast promise for every business, every society, and for all of you individually. & \textbf{01}\quad 它叫作区块链。区块链。这个词听起来并不是世界上最悦耳的词，但我相信，它现在就是下一代互联网，并且它对每一个企业、每一个社会，以及你们每一个人，都有巨大的前景。 \\ \hline
\textbf{02}\quad You know, for the past few decades we've had the internet of information, and when I send you an email or a PowerPoint file or something, I'm actually not sending you the original. I'm sending you a copy, and that's great. This is democratized information. But when it comes to assets, things like money, financial assets like stocks and bonds, loyalty points, intellectual property, music, art, a vote, carbon credit, and other assets, sending you a copy is a really bad idea. If I send you \$100, it's really important that I don't still have the money and that I can't send the same money to someone else. & \textbf{02}\quad 你知道，过去几十年里，我们拥有的是信息互联网。当我给你发一封邮件、一个 PowerPoint 文件或别的东西时，我实际上并没有把原件发给你。我发给你的是一份副本，而这很好。这使信息实现了民主化。但如果涉及资产，事情就不同了。钱、股票和债券这类金融资产、积分、知识产权、音乐、艺术品、一张选票、碳信用以及其他资产，如果只是发给你一份副本，那就是一个非常糟糕的想法。如果我给你 100 美元，非常重要的一点是：我不能还同时拥有这笔钱，也不能再把它发给你。 \\ \hline
\textbf{03}\quad So this has been called the double-spend problem by cryptographers for a long time. So today we rely entirely on big intermediaries, middlemen like banks, governments, big social media companies, credit card companies, and so on, to establish trust in our economy. And these intermediaries perform all the business and transaction logic of every kind of commerce, from authentication, identification of people through to clearing, settling, and record-keeping. And overall they do a pretty good job, but there are growing problems. & \textbf{03}\quad 密码学家很久以来把这称为“双重支付问题”。因此今天，我们完全依赖大型中介，也就是银行、政府、大型社交媒体公司、信用卡公司等中间人，来在我们的经济中建立信任。这些中介执行各种商业和交易逻辑，从认证、人员身份识别，一直到清算、结算和记录保存。总体来说，它们做得还不错，但问题正在增加。 \\ \hline
\textbf{04}\quad To begin, they're centralized. That means that they can be hacked and increasingly are. And J.P. Morgan, U.S. Federal Government, LinkedIn, Home Depot, and others found that out the hard way. They exclude billions of people from the global economy. For example, people who don't have enough money to have a bank account. They slow things down. It can take a second for an email to go around the world, but it can take days or weeks for money to move through the banking system across the city, and they take a big piece of the action, 10\% to 20\%, just to send money to another country. They capture our data, and that means that we can't monetize it or use it to better manage our lives, and our privacy is being undermined. And the biggest problem is that overall they've appropriated the largess of the digital age asymmetrically. And we have wealth creation, but we have growing social inequality. & \textbf{04}\quad 首先，它们是中心化的。这意味着它们可能被黑客攻击，而且这种情况越来越多。摩根大通、美国联邦政府、LinkedIn、Home Depot 等机构都已经痛苦地认识到了这一点。它们把数十亿人排除在全球经济之外，例如那些没有足够的钱去开银行账户的人。它们还让事情变慢。电子邮件一秒钟就可以传遍全球，但钱要通过银行系统在同一座城市里转移，可能需要几天甚至几周。它们还会从交易中拿走很大一块，比如仅仅向另一个国家汇钱就要收取 10\% 到 20\%。它们获取我们的数据，这意味着我们不能把这些数据变现，也不能用它更好地管理生活，而我们的隐私正在被削弱。最大的问题是，总体而言，它们不对称地占有了数字时代的丰厚成果。于是我们有了财富创造，但社会不平等也在增长。 \\ \hline
\end{longtable}


\clearpage
\section{Session 2.1 Activity 8: How Blockchain Works (Part 2)}
\addcontentsline{toc}{subsection}{中文译名：区块链如何运作（下）}

\renewcommand{\arraystretch}{1.23}
\begin{longtable}{>{\TranscriptColumn}p{0.47\textwidth} >{\TranscriptColumn}p{0.47\textwidth}}
\rowcolor{headerBg}\textbf{English Transcript} & \textbf{中文翻译} \\ \hline
\textbf{01}\quad So what if there were not only an internet of information, what if there were an internet of value? Some kind of vast, global, distributed ledger running on millions of computers and available to everybody, and where every kind of asset, from money to music, could be stored, moved, transacted, exchanged and managed, all without powerful intermediaries. What if there were a native medium for value? & \textbf{01}\quad 那么，如果不只有信息互联网，还存在一个价值互联网，会怎样？想象某种庞大的、全球性的、分布式的账本，它运行在数百万台计算机上，并且所有人都可以使用。在那里，从金钱到音乐的各种资产，都可以被存储、移动、交易、交换和管理，而不需要强大的中介机构。如果价值本身也有一种原生媒介，会怎样？ \\ \hline
\textbf{02}\quad Well, in 2008, the financial industry crashed, and perhaps propitiously, an anonymous person or persons named Satoshi Nakamoto created a paper where he developed a protocol for a digital cash that used an underlying cryptocurrency called Bitcoin. And this cryptocurrency enabled people to establish trust and do transactions without a third party. And the seemingly simple act set off a spark that's ignited the world, that has everyone excited or terrified or otherwise interested in many places. & \textbf{02}\quad 在 2008 年，金融业崩溃了。也许正是在这个恰当的时刻，一个名叫中本聪的匿名个人或群体写了一篇论文，在其中提出了一种数字现金协议，而它使用的底层加密货币叫作比特币。这种加密货币使人们能够在没有第三方的情况下建立信任并完成交易。这个看似简单的举动点燃了一粒火花，进而点燃了世界，让很多地方的人们感到兴奋、恐惧，或以其他方式产生兴趣。 \\ \hline
\textbf{03}\quad Now don't be confused about Bitcoin. Bitcoin is an asset. It goes up and down. And that should be of interest to you if you're a speculator. More broadly, it's a cryptocurrency. It's not a fiat currency controlled by a nation-state. And that's of greater interest. But the real pony here is the underlying technology. It's called blockchain. & \textbf{03}\quad 现在，不要把比特币搞混了。比特币是一种资产。它会上涨，也会下跌。如果你是投机者，这当然会让你感兴趣。更广泛地说，它是一种加密货币。它不是由某个民族国家控制的法定货币。而这更值得关注。但这里真正重要的东西，是它的底层技术。它叫作区块链。 \\ \hline
\textbf{04}\quad So for the first time now in human history, people everywhere can trust each other and transact, peer to peer. And trust is established not by some big institution, but by collaboration, by cryptography, and by some clever code. And because trust is native to the technology, I call this the trust protocol. & \textbf{04}\quad 所以，在人类历史上，现在第一次，世界各地的人可以彼此信任，并且点对点地交易。信任不是由某个大型机构建立的，而是由协作、密码学和一些巧妙的代码建立的。因为信任是这种技术原生的一部分，所以我把它称为“信任协议”。 \\ \hline
\textbf{05}\quad Now you're probably wondering, well, how does this thing work? Fair enough. Assets, digital assets like money to music and everything in between, are not stored in a central place. But they're distributed across a global ledger using the highest level of cryptography. And when a transaction is conducted, it's posted globally across millions and millions of computers. And out there around the world is a group of people called miners. These are not young people. They're Bitcoin miners. And they have massive computing power at their fingertips, 10 to 100 times bigger than all of Google worldwide. And these miners do a lot of work. And every 10 minutes, kind of like the heartbeat of a network, a block gets created that has all the transactions from the previous 10 minutes. And then the miners get to work trying to solve some tough problems. And they compete. And the first miner to find out the truth and to validate the block is rewarded in digital currency, in the case of the Bitcoin blockchain with Bitcoin. & \textbf{05}\quad 现在你可能会问：这个东西到底怎么运作？这个问题很合理。资产，也就是从金钱到音乐以及其间所有东西的数字资产，并不存放在一个中心地点。它们是用最高级别的密码学分布在一个全球账本上。当一笔交易发生时，它会被发布到全球数百万台计算机上。世界各地有一群人被称为矿工。他们不是年轻人，他们是比特币矿工。他们手边拥有巨大的计算能力，是全球 Google 总算力的 10 到 100 倍。这些矿工要做大量工作。每隔 10 分钟，就像网络的心跳一样，会产生一个区块，里面包含前 10 分钟的所有交易。然后矿工开始工作，试图解决一些很难的问题。他们相互竞争。第一个找出真相并验证区块的矿工会获得数字货币奖励；在比特币区块链中，奖励就是比特币。 \\ \hline
\textbf{06}\quad And then, this is the key part, that block is linked to the previous block and the previous block to create a chain of blocks. And every block is time-stamped, kind of like with a digital wax seal. So if I wanted to go in and hack a block and say, pay you and you with the same money, I'd have to hack that block, plus all the preceding blocks, the entire history of commerce on that blockchain, not just on one computer, but across millions of computers, simultaneously, all using the highest levels of encryption in the light of the most powerful computing resource in the world that's watching me, tough to do. And this is infinitely more secure than the computer systems that we have today. Blockchain, that's how it works. & \textbf{06}\quad 然后，关键部分来了：这个区块会连接到前一个区块，前一个区块再继续连接，从而形成一条区块链。每一个区块都有时间戳，有点像盖上了数字火漆印章。所以，如果我想进入系统篡改一个区块，比如用同一笔钱同时付给你和另一个人，我就必须篡改那个区块，还要篡改它之前所有的区块，也就是这条区块链上的整个商业历史。而且不是只在一台计算机上改，而是要同时在数百万台计算机上改，还要面对世界上最强大的计算资源在监视我，并且所有这些都使用最高等级的加密。要做到这一点非常困难。因此，它比我们今天拥有的计算机系统安全得多。区块链就是这样运作的。 \\ \hline
\end{longtable}


\clearpage
\section{Session 2.2 Activity 4: Automation in Workplaces and Its Implication (Part 1)}
\addcontentsline{toc}{subsection}{中文译名：职场自动化及其影响（上）}

\renewcommand{\arraystretch}{1.23}
\begin{longtable}{>{\TranscriptColumn}p{0.47\textwidth} >{\TranscriptColumn}p{0.47\textwidth}}
\rowcolor{headerBg}\textbf{English Transcript} & \textbf{中文翻译} \\ \hline
\textbf{01}\quad Delivering the future. Little robots taking parcels and groceries to the front door due to start in Greenwich, London next year. Up until now, the unpredictable complexity of the real world, infinite variety of hazards and obstacles meant the last mile of delivery needed a human. Not, it seems, any more. & \textbf{01}\quad 递送未来。小型机器人把包裹和食品杂货送到家门口，这项服务明年将在伦敦格林尼治启动。直到现在，现实世界不可预测的复杂性，以及无穷无尽的危险和障碍，意味着“最后一公里”配送需要人来完成。但现在看来，似乎不再需要了。 \\ \hline
\textbf{02}\quad Think of this. It's us waiting around at home in a five-hour delivery window. We have to take a day off work. It's so time-consuming and such a waste. Robotics and technology in the current day right now can solve those problems. So that's what we're doing. & \textbf{02}\quad 非常感谢。过去是我们在家里等着五小时的配送时间窗。我们不得不请一天假。这非常耗时，也非常浪费。而今天的机器人技术和相关技术现在就能解决这些问题。所以，这正是我们在做的事。 \\ \hline
\textbf{03}\quad There's no doubt that this coming wave of automation, of artificial intelligence and robots like this one, is going to deliver huge advantages to society. It's going to transform our lives, make them more convenient. The question is, is it also going to bring us some problems that society is going to struggle to cope with, not least, more inequality? & \textbf{03}\quad 毫无疑问，这一波即将到来的自动化、人工智能和像这样的机器人，会给社会带来巨大好处。它会改变我们的生活，让生活更便利。问题是，它是否也会带来一些社会难以应对的问题，尤其是更多的不平等？ \\ \hline
\textbf{04}\quad This is what could happen. With delivery bots, driverless vans, warehouse robots and online retailing, it's possible for a tiny number of people to sew up an entire market that once employed hundreds of thousands or even millions. The few can get fabulously rich. & \textbf{04}\quad 可能发生的是这样。通过配送机器人、无人驾驶货车、仓库机器人和在线零售，极少数人可能控制一个过去雇用了几十万甚至几百万人的整个市场。少数人会变得极其富有。 \\ \hline
\textbf{05}\quad As for the rest, in previous eras of job destruction, when mechanization chased farm workers off the land, it also created new and better-paid jobs in the towns and cities. And when muscle work disappeared, we moved into brain work. But soon, the machines will not only be stronger than us, but smarter too. Some new jobs will be created. These are robot handlers ready to take control if one of the delivery robots gets into trouble. The real worry is about how many of these new jobs there are going to be. & \textbf{05}\quad 至于其他人，在过去那些就业被破坏的时代，机械化把农场工人赶离土地，但同时也在城镇和城市里创造了新的、报酬更高的工作。当体力劳动消失时，我们转向了脑力劳动。但很快，机器不仅会比我们更强壮，也会比我们更聪明。一些新工作会被创造出来。这里的机器人操作员随时准备在某个配送机器人遇到麻烦时接管控制。真正令人担心的是：这种新工作到底会有多少？ \\ \hline
\textbf{06}\quad One study estimates that 35\% of UK jobs will be automated away in 10 to 20 years' time, and that jobs that pay under £30,000 a year are five times more at risk from automation than jobs that pay above £100,000. Research by McKinsey has warned that this transformation is happening 10 times faster and at 300 times the scale of the Industrial Revolution. Or, in other words, with 3,000 times the impact of the changes that transformed the Victorian world. These sculptures are perhaps a good metaphor for how most politicians have reacted to the challenge until now. & \textbf{06}\quad 一项研究估计，在未来 10 到 20 年内，英国 35\% 的工作将被自动化取代；年薪低于 3 万英镑的工作受到自动化威胁的风险，是年薪高于 10 万英镑工作的五倍。麦肯锡的研究警告说，这场转变的发生速度是工业革命的 10 倍，规模是工业革命的 300 倍。换句话说，它的影响将是那场改变维多利亚时代世界的变革的 3000 倍。这些雕塑也许很好地比喻了到目前为止大多数政治人物面对这一挑战时的反应。 \\ \hline
\end{longtable}


\clearpage
\section{Session 2.2 Activity 5: Automation in Workplaces and Its Implication (Part 2)}
\addcontentsline{toc}{subsection}{中文译名：职场自动化及其影响（下）}

\renewcommand{\arraystretch}{1.23}
\begin{longtable}{>{\TranscriptColumn}p{0.47\textwidth} >{\TranscriptColumn}p{0.47\textwidth}}
\rowcolor{headerBg}\textbf{English Transcript} & \textbf{中文翻译} \\ \hline
\textbf{01}\quad One idea they're going to be debating here in the European Parliament next month is this. If job-destroying robots are replacing tax-paying humans, well, perhaps we need to tax the robots. My proposal is that we monitor exactly what is happening on the job market, on the employment sector. And if robots are taking over more jobs than new jobs are created, we will be, or most member states and governance bodies, we'll be in a difficult position to collect enough money to finance all the services that we need. So if this is the case, I think we should think about taxing robots because someone has to pay for our normal life, for the infrastructures, for the services that the government should provide to the people. & \textbf{01}\quad 下个月他们将在欧洲议会辩论的一个想法是这样的：如果摧毁工作的机器人正在取代纳税的人类，那么也许我们需要向机器人征税。我的提议是，我们要准确监测就业市场、就业部门正在发生什么。如果机器人接管的工作多于新创造的工作，那么我们，或者说大多数成员国和治理机构，将会处在一个困难的位置，难以筹集足够的钱来资助我们所需要的所有服务。所以，如果情况确实如此，我认为我们应该考虑向机器人征税，因为总得有人为我们的正常生活、基础设施以及政府应该向人民提供的服务付钱。 \\ \hline
\textbf{02}\quad But how simple would this be? Some robots are easy to identify. The company replacing workers with a robot like Baxter here would be relatively straightforward to tax. But this wave of innovation is not all about the physical world. Much of this revolution, though, is taking place out of sight, inside reassuringly familiar-looking computer cases and server racks. But these new machines are learning and thinking. And learning and thinking is currently how much of humanity earns its living. & \textbf{02}\quad 但这会有多简单呢？有些机器人很容易识别。比如一家公司用这里的 Baxter 这样的机器人取代工人，对它征税会相对直接。但这一波创新并不全是关于物理世界的。事实上，这场革命有很大一部分发生在视线之外，发生在那些看上去令人熟悉、让人安心的电脑机箱和服务器机架内部。而这些新机器正在学习和思考。学习和思考，正是当前许多人类赖以谋生的方式。 \\ \hline
\textbf{03}\quad Welcome to SMAC. It stands for smart accountancy. This Berlin-based start-up has harnessed artificial intelligence in an attempt to revolutionise the boring old accountancy profession. The process looks familiar enough, clients of the service will scan their invoices and receipts into a computer. But inside, the software is doing something extraordinary. Instead of human accountants looking at the documents, the computer is figuring out the figures, adding them to the correct accounts. & \textbf{03}\quad 欢迎来到 SMAC。它代表“智能会计”。这家位于柏林的初创公司利用人工智能，试图彻底改变老旧而枯燥的会计行业。流程看起来很熟悉：这项服务的客户会把发票和收据扫描进电脑。但在内部，软件正在做一件非同寻常的事。不是由人类会计查看这些文件，而是由计算机弄清这些数字，并把它们加入正确的账户。 \\ \hline
\textbf{04}\quad Text on paper is pretty easy, but then you have to interpret the text, and the system has to know what is that number here and what is that number here, and where is the text information on the receipt. And that's interpreting the information, and that's really, in many, many areas, a revolution that we are seeing. For example, self-driving cars. Is that a person crossing the street, or is it just water on the street mirroring something? Can I go on driving, or do I have to stop? That's understanding, interpretation work. And that's also the work that has to be done here for the accounting purposes. The system really needs to do the interpretation. & \textbf{04}\quad 纸上的文字相当容易处理，但接下来你必须解释这些文字。系统必须知道这里的数字是什么、那里的数字是什么，文字信息在哪里，以及它意味着什么。这就是在解释信息，而在许许多多领域里，这正是我们正在看到的一场革命。比如自动驾驶汽车：那是一个正在过街的人，还是街上的积水在反射某个东西？我可以继续开车，还是必须停车？这就是理解、解释的工作。而为了会计目的，这里也必须完成同样的工作。系统确实需要进行解释。 \\ \hline
\textbf{05}\quad At SMAC, they say this will lead to less mundane, more interesting work for human accountants. But again, the question is, how many of these human jobs will survive? We are progressing very, very quickly in terms of innovation, but we have no idea where we are heading and where we want to reach. And I feel that at present, both governments and the private sector are quite happy to rely on neoliberalism and market forces and allow them to decide the direction of innovation. And this is extremely dangerous. & \textbf{05}\quad 在 SMAC，他们说这会让人类会计从乏味的工作转向更有意思的工作。但问题又来了：这些人类工作到底会有多少能保留下来？我们在创新方面进展非常非常快，但我们不知道自己正走向哪里，也不知道自己想要抵达哪里。我感觉目前政府和私营部门都相当乐意依赖新自由主义和市场力量，让市场决定创新的方向。而这极其危险。 \\ \hline
\textbf{06}\quad How dangerous? Well, don't take it from me. Want to know where we could be heading as a species? Have a look at our friend Oliver here. Just a century ago, millions of horses in the UK could earn their keep in agriculture, in the mines, in industry, in transport. The fact is that for almost every job that horses once monopolised, today they're not worth their feed and stabling. They couldn't give their labour away at any price. We can only hope that the robots are a little more sympathetic to us than we were to the horses. & \textbf{06}\quad 有多危险？好吧，不要只听我说。想知道作为一个物种，我们可能走向哪里吗？看看我们的朋友 Oliver。仅仅一个世纪前，英国有数百万匹马可以在农业、矿山、工业和运输业中养活自己。事实是，马曾经垄断的几乎每一份工作，今天都已经不值得为它们提供饲料和马厩了。它们即使以任何价格免费出卖劳动力，也没有人要。我们只能希望，机器人对我们能比我们当年对马更有同情心。 \\ \hline
\end{longtable}


\clearpage
\section{Session 2.2 Activity 7: The Digital Divide (Part 1)}
\addcontentsline{toc}{subsection}{中文译名：数字鸿沟（上）}

\renewcommand{\arraystretch}{1.23}
\begin{longtable}{>{\TranscriptColumn}p{0.47\textwidth} >{\TranscriptColumn}p{0.47\textwidth}}
\rowcolor{headerBg}\textbf{English Transcript} & \textbf{中文翻译} \\ \hline
\textbf{01}\quad Imagine this, you're sitting in a cozy cafe, sipping coffee, scrolling through your phone, watching this video. You have high-speed internet, access to unlimited information, and the ability to connect with people worldwide. Now, picture someone on the other side of the world, or maybe just in the next town over, who doesn't have this luxury. They don't have stable internet, a functioning device, or even the digital literacy to navigate today's technology-driven world. & \textbf{01}\quad 想象一下：你正坐在一家舒适的咖啡馆里，喝着咖啡，一边刷手机一边看这个视频。你有高速互联网，可以获取无限的信息，并且能够与世界各地的人联系。现在，想象世界另一端的人，或者也许只是隔壁城镇里的某个人，并没有这种便利。他们没有稳定的互联网，没有能正常使用的设备，甚至没有在今天这个由技术驱动的世界中行动所需的数字素养。 \\ \hline
\textbf{02}\quad Here's the question, in an age where everything is online, education, jobs, healthcare, and even government services, what happens to those who are left behind? Stay with me until the end, because I'm about to break down how the digital divide isn't just about having or not having internet access. It's about opportunity, power, and how technology is silently shaping a new kind of inequality. & \textbf{02}\quad 问题是，在一个一切都在线上的时代，教育、工作、医疗甚至政府服务都转到网上，那些被落下的人会怎样？请一直看到最后，因为我将要解释：数字鸿沟并不只是有无互联网接入的问题。它关乎机会、权力，以及技术如何在无声中塑造一种新的不平等。 \\ \hline
\textbf{03}\quad Technology was supposed to be the great equalizer. The internet promised a world where knowledge was free, opportunities were limitless, and anyone, regardless of background, could build a better future. But in reality, the more connected our world becomes, the deeper the gaps grow between those who have access and those who don't. The digital divide is about more than just owning a smartphone or a laptop. It's about who gets to participate in the digital economy and who gets left behind. & \textbf{03}\quad 技术本应成为伟大的均衡器。互联网曾承诺一个世界：知识是免费的，机会是无限的，任何人无论背景如何，都能建立更好的未来。但现实是，我们的世界越互联，拥有接入条件的人和没有接入条件的人之间的差距就越深。数字鸿沟不仅仅是拥有智能手机或笔记本电脑的问题。它关乎谁能参与数字经济，谁会被抛在后面。 \\ \hline
\textbf{04}\quad At its core, the digital divide exists at three levels. First, there's access, who has internet and devices, and who doesn't. Second, there's usability, who knows how to use technology effectively and who struggles with digital literacy. Third, there's empowerment, who benefits from technology, and who remains at a disadvantage, even when access exists. & \textbf{04}\quad 从核心来看，数字鸿沟存在于三个层面。第一是接入：谁有互联网和设备，谁没有。第二是可用性：谁知道如何有效使用技术，谁在数字素养上遇到困难。第三是赋权：谁从技术中受益，谁即使已经有了接入条件，仍然处于不利地位。 \\ \hline
\textbf{05}\quad Let's start with access. Over 2.6 billion people worldwide don't have internet access. That's nearly one third of the global population cut off from the modern world. And even among those who are online, the quality of access varies drastically. Some people enjoy high-speed fiber optic internet, while others are stuck with slow, unreliable connections that make even basic online tasks frustrating. & \textbf{05}\quad 让我们从接入开始。全球超过 26 亿人无法访问互联网。这接近全球人口的三分之一，他们被现代世界隔绝在外。即使在已经上线的人群中，接入质量也差异巨大。有些人享受高速光纤互联网，而另一些人却被缓慢、不可靠的连接困住，连基本的在线任务都令人沮丧。 \\ \hline
\textbf{06}\quad Think about education. During the COVID-19 pandemic, millions of students shifted to online learning. But what about the children in rural areas or low-income households who didn't have a laptop or reliable internet? They fell behind, not because they lacked intelligence or motivation, but because they lacked access. A lost year of education can mean fewer opportunities, lower income potential, and a future defined by limitations. & \textbf{06}\quad 想想教育。在新冠疫情期间，数百万学生转向线上学习。但那些农村地区或低收入家庭的孩子，如果没有笔记本电脑或可靠互联网怎么办？他们落后了，不是因为缺乏智力或动力，而是因为缺乏接入条件。失去一年的教育可能意味着更少的机会、更低的收入潜力，以及一个被限制所定义的未来。 \\ \hline
\textbf{07}\quad Then there's usability, the second layer of the digital divide. Even when people have internet access, not everyone knows how to use it effectively. Imagine a farmer in a developing country who has a smartphone but doesn't know how to access online markets to sell crops at a fair price. Or an elderly person who struggles to use online banking and gets scammed because they don't understand digital security. It's not just about having a device, it's about having the skills to use it in a way that improves life. & \textbf{07}\quad 然后是可用性，也就是数字鸿沟的第二层。即使人们可以上网，也不是每个人都知道如何有效使用它。想象一位发展中国家的农民有智能手机，却不知道如何进入线上市场，以公平价格出售农作物。或者一位老年人难以使用网上银行，并且因为不了解数字安全而被骗。这不只是拥有设备的问题，而是拥有用设备改善生活的技能的问题。 \\ \hline
\textbf{08}\quad And that brings us to the third, and perhaps the most dangerous, layer of the digital divide, empowerment. Even when people have access and skills, the real question is, who benefits from technology? Big corporations, governments, and tech elites are leveraging digital advancements to expand their power, while marginalized communities often find themselves on the losing end. Job markets are shifting toward automation and digital skills. Those who can adapt thrive, while those who can't are pushed out of the workforce. & \textbf{08}\quad 这就把我们带到数字鸿沟的第三层，也许也是最危险的一层：赋权。即使人们拥有接入和技能，真正的问题仍然是：谁从技术中受益？大型企业、政府和科技精英正在利用数字进步扩大自己的权力，而边缘化社区却常常落后。就业市场正在转向自动化和数字技能。能够适应的人会兴旺，不能适应的人会被挤出劳动力市场。 \\ \hline
\end{longtable}


\clearpage
\section{Session 2.2 Activity 8: The Digital Divide (Part 2)}
\addcontentsline{toc}{subsection}{中文译名：数字鸿沟（下）}

\renewcommand{\arraystretch}{1.23}
\begin{longtable}{>{\TranscriptColumn}p{0.47\textwidth} >{\TranscriptColumn}p{0.47\textwidth}}
\rowcolor{headerBg}\textbf{English Transcript} & \textbf{中文翻译} \\ \hline
\textbf{01}\quad Think about this, tech jobs, freelancing and remote work opportunities are booming, but what if you don't have a laptop or stable internet? What if you've never learned how to code or navigate digital platforms? The divide is no longer just about who has technology, it's about who gets to use it to create a better life. & \textbf{01}\quad 想想这一点：技术工作、自由职业和远程工作机会正在蓬勃发展。但如果你没有笔记本电脑或稳定的互联网，会怎样？如果你从来没有学过编程，也不会使用数字平台，会怎样？如今的鸿沟不再只是关于谁拥有技术，而是关于谁能利用技术创造更好的生活。 \\ \hline
\textbf{02}\quad The digital divide isn't just a developing world issue. Even in rich countries, the gap is growing. In the US, millions of rural communities lack reliable broadband. In major cities, low-income neighborhoods often have fewer public Wi-Fi spots, outdated computer labs and limited digital literacy programs. & \textbf{02}\quad 数字鸿沟并不只是发展中世界的问题。即使在富裕国家，差距也在扩大。在美国，数百万农村社区缺乏可靠的宽带。在大城市，低收入社区往往公共 Wi-Fi 点更少，电脑实验室设备陈旧，数字素养培训也有限。 \\ \hline
\textbf{03}\quad And here's where it gets even more serious, that digital divide isn't just about convenience, it's about survival. During the pandemic, people with internet access could work remotely. Those without it? Many lost their jobs. Access to healthcare is also shifting online. Telemedicine, online health portals and digital prescriptions are becoming the norm. But what about those who can't afford internet or don't know how to use digital healthcare services? They receive worse medical care simply because they lack access to a screen. & \textbf{03}\quad 更严重的是，数字鸿沟并不只是便利问题，而是生存问题。疫情期间，拥有互联网接入的人可以远程工作。没有互联网的人呢？许多人失去了工作。医疗服务的获取也正在转向线上。远程医疗、在线健康门户和数字处方正在成为常态。但那些负担不起互联网，或不知道如何使用数字医疗服务的人怎么办？他们只是因为无法接触屏幕，就得到更差的医疗照护。 \\ \hline
\textbf{04}\quad And think about democracy. Political debates, news updates and even voting are increasingly digital. When people lack digital access, they're not just disconnected from technology, they're disconnected from decision-making. & \textbf{04}\quad 再想想民主。政治辩论、新闻更新，甚至投票都越来越数字化。当人们缺乏数字接入时，他们不只是与技术断开连接，也被剥夺了参与塑造未来的讨论的机会。 \\ \hline
\end{longtable}


\clearpage
\section{Session 2.3 Activity 7 \& 8: How to Keep Human Bias Out of AI}
\addcontentsline{toc}{subsection}{中文译名：如何避免人类偏见进入人工智能}

\renewcommand{\arraystretch}{1.23}
\begin{longtable}{>{\TranscriptColumn}p{0.47\textwidth} >{\TranscriptColumn}p{0.47\textwidth}}
\rowcolor{headerBg}\textbf{English Transcript} & \textbf{中文翻译} \\ \hline
\textbf{01}\quad AI is being used to help decide whether or not you get that job interview, how much you pay for your car insurance, how good your credit score is, and even what rating you get in your annual performance review. But these decisions are all being filtered through its assumptions about our identity, our race, our gender, our age. How is that happening? & \textbf{01}\quad AI 正在被用来帮助决定你是否能得到那次工作面试、你的车险要付多少钱、你的信用评分有多好，甚至你在年度绩效评估中会得到什么评级。但这些决定都被 AI 关于我们身份、种族、性别和年龄的假设过滤了一遍。这是怎么发生的？ \\ \hline
\textbf{02}\quad Now, imagine an AI is helping a hiring manager find the next tech leader in the company. So far, the manager has been hiring mostly men, so the AI learns men are more likely to be programmers than women. And it's a very short leap from there to the idea that men make better programmers than women. We have reinforced our own bias into the AI, and now it's screening out female candidates. & \textbf{02}\quad 现在，想象一个 AI 正在帮助招聘经理寻找公司下一位技术负责人。到目前为止，这位经理雇用的大多是男性，于是 AI 学到：男性比女性更有可能成为程序员。很快，它就会从这一点滑到“男性比女性更适合做程序员”。我们把自己的偏见强化进了 AI，现在它开始筛掉女性候选人。 \\ \hline
\textbf{03}\quad Hang on, if a human hiring manager did that, we'd be outraged. We wouldn't allow it. This kind of gender discrimination is not OK. And yet, somehow, AI has become above the law because the machine made the decision. That's not it. We are also reinforcing our bias in how we interact with AI. How often do you use a voice assistant like Siri, Alexa or even Cortana? They all have two things in common. One, they can never get my name right. And second, they are all female. They're designed to be your obedient servants, turning your lights on and off, ordering your shopping. You get male AIs too, but they tend to be more high-powered, like IBM Watson making business decisions, Salesforce Einstein or Ross, the robot lawyer. So, poor robots, even they suffer from sexism in the workplace. & \textbf{03}\quad 等等，如果一个人类招聘经理这么做，我们会愤怒。我们不会允许这种事。这种性别歧视是不可以的。然而，不知怎么地，因为是机器做出了决定，AI 就好像被置于法律之上了。不止如此。我们还在与 AI 互动的方式中强化自己的偏见。你多久使用一次像 Siri、Alexa 甚至 Cortana 这样的语音助手？它们都有两个共同点。第一，它们永远念不准我的名字。第二，它们都是女性化的。它们被设计成你顺从的仆人，帮你开灯关灯、帮你下购物订单。你也会看到男性化的 AI，但它们往往更高权力，比如做商业决策的 IBM Watson、Salesforce Einstein，或者机器人律师 Ross。所以，可怜的机器人，连它们也在工作场所遭受性别主义。 \\ \hline
\textbf{04}\quad Think about how these two things combine and affect a kid growing up in today's world around AI. So, they're doing some research for a school project, and they Google images of CEO. The algorithm shows them results of mostly men. And now, they Google personal assistant. As you can guess, it shows them mostly females. And then, they want to put on some music and maybe order some food. And now, they're barking orders at a female voice, an obedient female voice assistant. Some of our brightest minds are creating this technology today, technology that they could have created in any way they wanted. And yet, they've chosen to create it in the style of 1950s Mad Men secretary. Yay! & \textbf{04}\quad 想一想，这两件事结合起来，会怎样影响一个在今天这个 AI 世界中成长的孩子。比如他们正在为学校项目做研究，在 Google 图片里搜索 CEO。算法给他们显示的结果大多是男性。然后他们搜索 personal assistant。你可以猜到，结果显示的大多是女性。接着，他们想放点音乐，也许再点些食物。于是他们又开始向一个女性声音、一个顺从的女性语音助手发号施令。今天一些最聪明的人正在创造这种技术，而他们本可以用任何方式创造它。但他们却选择把它做成 20 世纪 50 年代《广告狂人》里秘书的风格。好耶。 \\ \hline
\textbf{05}\quad But okay, don't worry. This is not going to end with me telling you that we're all heading towards sexist, racist machines running the world. The good news about AI is that it is entirely within our control. We get to teach the right values, the right ethics to AI. So, there are three things we can do. One, we can be aware of our own biases and the bias in machines around us. Two, we can make sure there are diverse teams building this technology. And three, we have to give it diverse experiences to learn from. & \textbf{05}\quad 不过，好吧，别担心。我的结尾不是要告诉你：我们都正走向由性别歧视、种族歧视的机器统治世界。关于 AI 的好消息是，它完全在我们的控制之内。我们可以把正确的价值观、正确的伦理教给 AI。所以我们可以做三件事。第一，我们可以意识到自己的偏见，以及周围机器中的偏见。第二，我们可以确保由多元化团队来建设这项技术。第三，我们必须给它多样化的经验去学习。 \\ \hline
\textbf{06}\quad I can talk about the first two from personal experience. When you work in technology and you don't look like a Mark Zuckerberg or Elon Musk, your life is a little bit difficult, your ability gets questioned. Here's just one example. Like most developers, I often join online tech forums and share my knowledge to help others. And I found when I log on as myself with my own photo, my own name, I tend to get questions or comments like this. What makes you think you're qualified to talk about AI? What makes you think you know about machine learning? So, as you do, I made a new profile. And this time, instead of my own picture, I chose a cat with a jetpack on it. And I chose a name that did not reveal my gender. You can probably guess where this is going, right? So, this time, I didn't get any of those patronizing comments about my ability, and I was able to actually get some work done. & \textbf{06}\quad 我可以用个人经历来谈前两点。当你在技术领域工作，而你看起来不像 Mark Zuckerberg 或 Elon Musk 时，你的生活会有点困难，你的能力会受到质疑。这里只举一个例子。像大多数开发者一样，我经常加入在线技术论坛，分享知识来帮助别人。我发现，当我用自己的照片和自己的名字登录时，我往往会收到这样的问题或评论：“你凭什么觉得自己有资格谈 AI？”“你凭什么觉得你懂机器学习？”所以，像很多人会做的那样，我创建了一个新资料。这一次，我没有用自己的照片，而是选了一只背着喷气背包的猫。我还选了一个不会暴露我性别的名字。你大概能猜到接下来会怎样，对吧？这一次，我没有再收到那些居高临下质疑我能力的评论，我也终于能真正完成一些工作。 \\ \hline
\textbf{07}\quad And it sucks, guys. I've been building robots since I was 15, I have a few degrees in computer science, and yet, I had to hide my gender in order for my work to be taken seriously. So, what's going on here? Are men just better at technology than women? Another study found that when women coders on one platform hid their gender, like myself, their code was accepted 4\% more than men's code. So, this is not about the talent. This is about an elitism in AI that says a programmer needs to look like a certain person. & \textbf{07}\quad 而这真的很糟糕。伙计们，我从 15 岁起就在做机器人，我有几个计算机科学学位，可是为了让自己的工作被认真对待，我仍然不得不隐藏自己的性别。那么这里到底发生了什么？难道男性天生就比女性更擅长技术吗？另一项研究发现，在一个平台上，当女性程序员像我一样隐藏性别时，她们的代码被接受的比例比男性高 4\%。所以这不是能力问题。这是 AI 领域的一种精英主义：它认为程序员必须长得像某一种人。 \\ \hline
\end{longtable}


\clearpage
\section{Session 2.3 Activity 14: Ethics of AI: Challenges and Governance}
\addcontentsline{toc}{subsection}{中文译名：人工智能伦理：挑战与治理}

\renewcommand{\arraystretch}{1.23}
\begin{longtable}{>{\TranscriptColumn}p{0.47\textwidth} >{\TranscriptColumn}p{0.47\textwidth}}
\rowcolor{headerBg}\textbf{English Transcript} & \textbf{中文翻译} \\ \hline
\textbf{01}\quad Do you use any navigation application to go through traffic jams? What happens when you scroll down through your social media feed? Do you follow any recommendation that any streaming service might give you? So by the end of the day, we know that AI is there, but do we actually understand what is going on, and can we trust the output of those applications? & \textbf{01}\quad 你会使用导航应用来穿过交通拥堵吗？当你向下滑动社交媒体信息流时，发生了什么？你会听从某个流媒体服务给你的任何推荐吗？所以，到一天结束时，我们知道 AI 就在那里，但我们真的理解正在发生什么吗？我们能信任这些应用的输出吗？ \\ \hline
\textbf{02}\quad To date, the approach that we've taken to these kinds of technologies is to say, well, consumers can figure this out on their own. They can read the terms and conditions on websites. They can choose not to participate in certain digital environments, but increasingly, these products and platforms are part of our lives. They are part of the way that we provide education. They are part of the way that we look for jobs. And that power imbalance cannot be addressed by simply giving consumers more information or giving them individual rights of complaint. & \textbf{02}\quad 到目前为止，我们对这类技术采取的做法是说：消费者可以自己弄明白。他们可以阅读网站上的条款和条件。他们可以选择不参与某些数字环境。但越来越多地，这些产品和平台已经成为我们生活的一部分。它们是我们提供教育的方式的一部分。它们是我们寻找工作的方式的一部分。而这种权力不平衡，不能只靠给消费者更多信息，或者给他们个人投诉权来解决。 \\ \hline
\textbf{03}\quad What we're ultimately going to have to do if we're going to change the kind of structural ways in which these technologies are designed is to push the responsibilities back on our designers, back on the organizations who are relying on technologies to change their practices. Technologies like artificial intelligence have the potential to empower people and to broaden their perspectives, or they can widen inequalities and really don't contribute to address our societal challenges. But we cannot blame the technologies. It's not about the technologies. It's about us developing the framework, shaping the rules that will allow these technologies to deliver to what we want them to achieve. & \textbf{03}\quad 如果我们要改变这些技术在结构层面被设计的方式，我们最终必须把责任推回给设计者，推回给那些部署这些技术的组织，要求它们改变实践。像人工智能这样的技术有潜力赋能人们、拓宽他们的视野，也可能扩大不平等，并让我们无法应对社会挑战。但我们不能责怪技术本身。问题不在技术本身，而在于我们要建立框架、塑造规则，使这些技术能够实现我们希望它们实现的目标。 \\ \hline
\textbf{04}\quad And we need to make sure that we embed those very important ethical principles, the protection and promotion of human rights and human dignity that will at the end determine the outcomes. Responsible governance of AI technologies without the buy-in of Big Tech and other companies is impossible. It is necessary for us to find ways in which to explain to them that ethics is not about abstract principles. Ethics should be bottom up. It is a dynamic system. It should enable innovation and in the end lead to trust in products of these companies which would lead to success in terms of whatever their business aims are. & \textbf{04}\quad 我们需要确保把非常重要的伦理原则嵌入其中，也就是保护和促进人权与人的尊严；最终正是这些原则会决定结果。没有大型科技公司和其他公司的参与，AI 技术的负责任治理是不可能的。我们必须找到方式向它们解释：伦理不是抽象原则。伦理应当自下而上。它是一个动态系统。它应当促成创新，并最终带来对这些公司产品的信任，而这种信任又会导向它们商业目标上的成功。 \\ \hline
\textbf{05}\quad We see how ethical debates end up playing an important role in shaping the conversation of AI regulation over the last five years. We saw a plethora of charters, declarations of AI ethical principles. We are seeing right now that those principles are being applied on a very practical basis. A lot of countries in Latin America have already come up with their national strategies on artificial intelligence. But some countries in the region are going a step further. They are actually regulating through hard law some of the principles of artificial intelligence. In the European Union, we're looking at draft AI acts. In the US, Congress is starting to talk in great detail about monopoly power of technology companies. And you can see the wheels turning across the world as countries move from awareness-raising to strategy, to attempts at implementation, attempts at regulation. & \textbf{05}\quad 我们看到，在过去五年里，伦理辩论在塑造 AI 监管的讨论方面发挥了重要作用。我们看到了大量 AI 伦理原则的宪章和宣言。现在我们也看到，这些原则正在非常实际的基础上被应用。拉丁美洲很多国家已经提出了自己的人工智能国家战略。但该地区有些国家还更进一步，正在通过硬法来监管人工智能的一些原则。在欧盟，我们正在看 AI 法案草案。在美国，国会开始非常详细地讨论科技公司的垄断权力。你可以看到，世界各地的齿轮正在转动，各国正在从提高意识，走向战略，再走向实施尝试和监管尝试。 \\ \hline
\end{longtable}


\clearpage
\section{Session 2.4 Activity 4: How I'm Fighting Bias in Algorithms (Part 1)}
\addcontentsline{toc}{subsection}{中文译名：我如何对抗算法偏见（上）}

\renewcommand{\arraystretch}{1.23}
\begin{longtable}{>{\TranscriptColumn}p{0.47\textwidth} >{\TranscriptColumn}p{0.47\textwidth}}
\rowcolor{headerBg}\textbf{English Transcript} & \textbf{中文翻译} \\ \hline
\textbf{01}\quad So how did this happen? Why am I sitting in front of a computer in a white mask trying to be detected by a cheap webcam? Well, when I'm not fighting the coded gaze as a poet of code, I'm a graduate student at the MIT Media Lab. And there I have the opportunity to work on all sorts of whimsical projects, including the Aspire Mirror, a project I did so I could project digital masks onto my reflection so in the morning, if I wanted to feel powerful, I could put on a lion mask. If I wanted to be uplifted, I might have a quote. & \textbf{01}\quad 那么，这是怎么发生的？为什么我会坐在电脑前，戴着一个白色面具，试图让一个廉价摄像头检测到我？当我不是作为代码诗人去对抗“编码凝视”时，我是 MIT 媒体实验室的一名研究生。在那里，我有机会做各种充满奇思妙想的项目，包括 Aspire Mirror。这个项目让我可以把数字面具投射到我在镜子里的倒影上。如果早上我想感觉强大，我可以投射一张狮子面具。如果我想受到鼓舞，我可以显示一句引语。 \\ \hline
\textbf{02}\quad So I used generic facial recognition software to build the system, but found that it was really hard to test it unless I wore a white mask. Unfortunately, I've run into this issue before. When I was an undergraduate at Georgia Tech studying computer science, I used to work on social robots. And one of my tasks was to get a robot to play peek-a-boo, a simple turn-taking game where partners covered their face and then uncover it, saying peek-a-boo. The problem is peek-a-boo doesn't really work if I can't see you, and my robot couldn't see me. & \textbf{02}\quad 于是我用通用人脸识别软件来搭建这个系统，但我发现，除非我戴上白色面具，否则测试起来非常困难。不幸的是，我以前也遇到过这个问题。当我还在佐治亚理工读计算机科学本科时，我曾经做社交机器人。其中一项任务是让机器人玩躲猫猫，这是一个简单的轮流游戏：伙伴们遮住脸，然后露出来，说“peek-a-boo”。问题是，如果我看不见你，躲猫猫就不能真正进行；而我的机器人看不见我。 \\ \hline
\textbf{03}\quad But I borrowed my roommate's face to get the project done, submitted the assignment, and figured, you know what? Somebody else will solve this problem. Not too long after, I was in Hong Kong for an entrepreneurship competition. The organizers decided to take participants on a tour of local startups. One of the startups had a social robot, and they decided to do a demo. The demo worked on everybody until it got to me, and you can probably guess it. It couldn't detect my face. & \textbf{03}\quad 于是我借用了室友的脸来完成项目，提交了作业，并心想：你知道吗，总会有别人解决这个问题的。没过多久，我去香港参加一次创业比赛。组织者决定带参赛者参观当地初创公司。其中一家初创公司有一个社交机器人，他们决定做一次演示。演示对每个人都有效，直到轮到我；你大概可以猜到结果。它检测不到我的脸。 \\ \hline
\textbf{04}\quad I asked the developers what was going on, and it turned out we had used the same generic facial recognition software. Halfway around the world, I learned that algorithmic bias can travel as quickly as it takes to download some files off of the internet. & \textbf{04}\quad 我问开发人员发生了什么，结果发现我们使用的是同一套通用人脸识别软件。跨越半个世界，我学到：算法偏见传播的速度，可以像从互联网上下载几个文件一样快。 \\ \hline
\end{longtable}


\clearpage
\section{Session 2.4 Activity 5: How I'm Fighting Bias in Algorithms (Part 2)}
\addcontentsline{toc}{subsection}{中文译名：我如何对抗算法偏见（下）}

\renewcommand{\arraystretch}{1.23}
\begin{longtable}{>{\TranscriptColumn}p{0.47\textwidth} >{\TranscriptColumn}p{0.47\textwidth}}
\rowcolor{headerBg}\textbf{English Transcript} & \textbf{中文翻译} \\ \hline
\textbf{01}\quad So what's going on? Why isn't my face being detected? Well, we have to look at how we give machine sight. Computer vision uses machine learning techniques to do facial recognition. So how this works is you create a training set with examples of faces. This is a face, this is a face, this is not a face. And over time, you can teach a computer how to recognize other faces. & \textbf{01}\quad 那么，发生了什么？为什么我的脸没有被检测到？我们必须看看我们是怎样赋予机器视觉的。计算机视觉使用机器学习技术来做人脸识别。它的工作方式是：你创建一个训练集，里面有一些人脸示例。这是一张脸，这是一张脸，这不是一张脸。随着时间推移，你就能教会计算机如何识别其他脸。 \\ \hline
\textbf{02}\quad However, if the training sets aren't really that diverse, any face that deviates too much from the established norm will be harder to detect, which is what was happening to me. But don't worry, there's some good news. Training sets don't just materialize out of nowhere. We actually can create them. So there's an opportunity to create full spectrum training sets that reflect a richer portrait of humanity. & \textbf{02}\quad 但是，如果训练集并不真正多样化，任何偏离既定常模太多的脸都会更难被检测到，而这正是发生在我身上的事。不过别担心，也有一些好消息。训练集不会凭空出现。我们实际上可以创造它们。因此，我们有机会创造全光谱训练集，让它们反映出更丰富的人类肖像。 \\ \hline
\textbf{03}\quad Now, you've seen it in my examples, how social robots were how I found out about exclusion through algorithmic bias. But algorithmic bias can also lead to discriminatory practices. Across the US, police departments are starting to use facial recognition software in their crime-fighting arsenal. Georgetown Law published a report showing that one in two adults in the US, that's 117 million people, have their faces in facial recognition networks. Police departments can currently look at these networks unregulated using algorithms that have not been audited for accuracy. & \textbf{03}\quad 现在，你已经在我的例子中看到，社交机器人让我发现了算法偏见导致的排斥。但算法偏见也可能导致歧视性做法。在美国，警察部门正在开始把人脸识别软件用于他们的打击犯罪工具库。乔治城法学院发表了一份报告，显示美国每两个成年人中就有一个，也就是 1.17 亿人，已经把自己的脸放进了人脸识别网络。警察部门目前可以在不受监管的情况下查看这些网络，而且使用的是没有经过准确性审计的算法。 \\ \hline
\textbf{04}\quad Yet we know facial recognition is not foolproof and labeling faces consistently remains a challenge. You might have seen this on Facebook. My friends and I laugh all the time when we see other people mislabeled in our photos. But misidentifying a suspected criminal is no laughing matter, nor is breaching civil liberties. & \textbf{04}\quad 然而我们知道，人脸识别并非万无一失，持续准确地给人脸打标签仍然是一个挑战。你可能在 Facebook 上见过这种情况。我和朋友们看到照片里其他人被错误标注时，总是会笑。但把犯罪嫌疑人识别错并不是笑话，侵犯公民自由也不是笑话。 \\ \hline
\textbf{05}\quad Machine learning is being used for facial recognition, but it's also extending beyond the realm of computer vision. In her book, Weapons of Math Destruction, data scientist Cathy O'Neil talks about the rising new WMDs. Widespread, mysterious and destructive algorithms that are increasingly being used to make decisions that impact more aspects of our lives. So who gets hired or fired? Do you get that loan? Do you get insurance? Are you admitted into the college that you wanted to get into? Do you and I pay the same price for the same product purchased on the same platform? & \textbf{05}\quad 机器学习正在被用于人脸识别，但它也正在扩展到计算机视觉之外。在她的书《数学毁灭性武器》中，数据科学家 Cathy O'Neil 谈到了正在兴起的新型“大规模杀伤性武器”：广泛、神秘且具有破坏性的算法，越来越多地被用来做出影响我们生活更多方面的决定。那么，谁会被雇用或解雇？你能不能拿到贷款？你能不能买到保险？你会不会被你想进的大学录取？你和我在同一平台购买同一产品时，会不会支付同样的价格？ \\ \hline
\textbf{06}\quad Law enforcement is also starting to use machine learning for predictive policing. Some judges use machine-generated risk scores to determine how long an individual is going to spend in prison. So we really have to think about these decisions. Are they fair? And we've seen that algorithmic bias doesn't necessarily always lead to fair outcomes. & \textbf{06}\quad 执法部门也开始使用机器学习进行预测性警务。一些法官使用机器生成的风险评分，来决定一个人将在监狱里待多久。因此，我们真的必须思考这些决定。它们公平吗？我们已经看到，算法偏见并不一定总是带来公平的结果。 \\ \hline
\end{longtable}


\clearpage
\section{Session 2.4 Activity 7: The Ethical Dilemma of Self-Driving Cars (Part 1)}
\addcontentsline{toc}{subsection}{中文译名：自动驾驶汽车的伦理困境（上）}

\renewcommand{\arraystretch}{1.23}
\begin{longtable}{>{\TranscriptColumn}p{0.47\textwidth} >{\TranscriptColumn}p{0.47\textwidth}}
\rowcolor{headerBg}\textbf{English Transcript} & \textbf{中文翻译} \\ \hline
\textbf{01}\quad This is a thought experiment. Let's say at some point in the not-so-distant future you're barreling down the highway in your self-driving car, and you find yourself boxed in on all sides by other cars. Suddenly, a large, heavy object falls off the truck in front of you. Your car can't stop in time to avoid the collision, so it needs to make a decision. Go straight and hit the object, swerve left into an SUV, or swerve right into a motorcycle. & \textbf{01}\quad 这是一个思想实验。假设在不远的未来某个时候，你正坐在自动驾驶汽车里沿高速公路疾驰，发现自己被四面八方的其他车辆围住。突然，前方卡车上掉下一个又大又重的物体。你的车来不及停下以避免碰撞，所以它必须做一个决定：直行撞上那个物体，向左转撞上一辆 SUV，还是向右转撞上一辆摩托车。 \\ \hline
\textbf{02}\quad Should it prioritize your safety by hitting the motorcycle? Minimize danger to others by not swerving, even if it means hitting the large object and sacrificing your life. Or take the middle ground by hitting the SUV, which has a high passenger safety rating. So what should the self-driving car do? & \textbf{02}\quad 它是否应该通过撞向摩托车来优先保护你的安全？是否应该为了把对他人的危险降到最低而不转向，即使这意味着撞上那个大物体并牺牲你的生命？或者它应该采取中间方案，撞向那辆乘客安全评级较高的 SUV？那么，自动驾驶汽车到底应该怎么做？ \\ \hline
\textbf{03}\quad If we were driving that boxed-in car in manual mode, whichever way we'd react would be understood as just that, a reaction, not a deliberate decision. It would be an instinctual panicked move with no forethought or malice. But if a programmer were to instruct the car to make the same move, given conditions it may sense in the future, well, that looks more like premeditated homicide. & \textbf{03}\quad 如果我们是在手动模式下驾驶那辆被困住的车，不管我们往哪个方向反应，都会被理解为反应，而不是一个深思熟虑的决定。那会是出于本能的恐慌动作，没有事先考虑，也没有恶意。但如果一名程序员指示汽车在未来感知到某些条件时做出同样的动作，那看起来就更像是有预谋的杀人。 \\ \hline
\textbf{04}\quad Now, to be fair, self-driving cars are predicted to dramatically reduce traffic accidents and fatalities by removing human error from the driving equation. Plus, there may be all sorts of other benefits - eased road congestion, decreased harmful emissions, and minimized unproductive and stressful driving time. But accidents can and will still happen, and when they do, their outcomes may be determined months or years in advance by programmers or policymakers. & \textbf{04}\quad 公平地说，自动驾驶汽车预计会通过消除驾驶中的人为错误，大幅减少交通事故和死亡人数。另外，它们也可能带来各种其他好处，比如缓解道路拥堵、减少有害排放，并减少低效且令人紧张的驾驶时间。但事故仍然可能发生，而且一定会发生；当事故发生时，其结果可能早在几个月甚至几年之前，就已经由程序员或政策制定者决定好了。 \\ \hline
\end{longtable}


\clearpage
\section{Session 2.4 Activity 8: The Ethical Dilemma of Self-Driving Cars (Part 2)}
\addcontentsline{toc}{subsection}{中文译名：自动驾驶汽车的伦理困境（下）}

\renewcommand{\arraystretch}{1.23}
\begin{longtable}{>{\TranscriptColumn}p{0.47\textwidth} >{\TranscriptColumn}p{0.47\textwidth}}
\rowcolor{headerBg}\textbf{English Transcript} & \textbf{中文翻译} \\ \hline
\textbf{01}\quad And they'll have some difficult decisions to make. It's tempting to offer up general decision-making principles like minimize harm, but even that quickly leads to morally murky decisions. For example, let's say we have the same initial setup, but now there's a motorcyclist wearing a helmet to your left and another one without a helmet to your right. Which one should your robot car crash into? & \textbf{01}\quad 它们会有一些困难的决定要做。提出“最小化伤害”这样的通用决策原则很诱人，但即使是这种原则，也会很快导向道德上模糊的决定。例如，假设我们有同样的初始设置，但现在你的左边有一名戴头盔的摩托车手，右边还有一名没有戴头盔的摩托车手。你的机器人汽车应该撞向哪一个？ \\ \hline
\textbf{02}\quad If you say the biker with the helmet because she's more likely to survive, then aren't you penalizing the responsible motorist? If instead you say the biker without the helmet because he's acting irresponsibly, then you've gone way beyond the initial design principle about minimizing harm, and the robot car is now meting out street justice. & \textbf{02}\quad 如果你说撞向戴头盔的骑手，因为她更可能活下来，那么你是不是在惩罚负责任的驾驶者？如果你反而说撞向没戴头盔的骑手，因为他行为不负责任，那么你就已经远远超出了最初关于“最小化伤害”的设计原则，机器人汽车成了在执行某种街头正义。 \\ \hline
\textbf{03}\quad The ethical considerations get more complicated here. In both of our scenarios, the underlying design is functioning as a targeting algorithm of sorts. In other words, it's systematically favoring or discriminating against a certain type of object to crash into. And the owners of the target vehicles will suffer the negative consequences of this algorithm through no fault of their own. & \textbf{03}\quad 这里的伦理考量变得更复杂。在我们的两个情景中，底层设计实际上像某种目标选择算法。换句话说，它系统性地偏向或歧视某一类将被撞击的对象。而目标车辆的车主，会在自己没有任何过错的情况下，承受这个算法带来的负面后果。 \\ \hline
\textbf{04}\quad Our new technologies are opening up many other novel ethical dilemmas. For instance, if you had to choose between a car that would always save as many lives as possible in an accident, or one that would save you at any cost, which would you buy? What happens if the cars start analyzing and factoring in the passengers of the cars and the particulars of their lives? Could it be the case that a random decision is still better than a predetermined one designed to minimize harm? And who should be making all of these decisions anyhow? Programmers? Companies? Governments? & \textbf{04}\quad 我们的新技术正在打开许多其他新的伦理困境。比如，如果你必须在两种车之间选择：一种车在事故中总是尽可能拯救最多人的生命，另一种车无论付出什么代价都救你，你会买哪一种？如果汽车开始分析并考虑车上乘客以及他们生活中的具体情况，会发生什么？有没有可能，随机决策仍然比一个预先设定、旨在最小化伤害的决定更好？而这些决定究竟应该由谁来做？程序员？公司？政府？ \\ \hline
\textbf{05}\quad Reality may not play out exactly like our thought experiments, but that's not the point. They're designed to isolate and stress test our intuitions on ethics, just like science experiments do for the physical world. Spotting these moral hairpin turns now will help us maneuver the unfamiliar road of technology ethics and allow us to cruise confidently and conscientiously into our brave new future. & \textbf{05}\quad 现实可能不会完全像我们的思想实验那样展开，但那不是重点。思想实验的设计目的，是分离并压力测试我们关于伦理的直觉，就像科学实验为物理世界所做的那样。现在发现这些道德急弯，将帮助我们在技术伦理这条陌生道路上操控方向，并让我们能够更自信、更有良知地驶向勇敢的新未来。 \\ \hline
\end{longtable}


\clearpage
\section{Session 2.4 Activity 14: Let's Not Use Mars as a Backup Planet}
\addcontentsline{toc}{subsection}{中文译名：不要把火星当作备用星球}

\renewcommand{\arraystretch}{1.23}
\begin{longtable}{>{\TranscriptColumn}p{0.47\textwidth} >{\TranscriptColumn}p{0.47\textwidth}}
\rowcolor{headerBg}\textbf{English Transcript} & \textbf{中文翻译} \\ \hline
\textbf{01}\quad We're at a tipping point in human history, a species poised between gaining the stars and losing the planet we call home. Even in just the past few years, we've greatly expanded our knowledge of how Earth fits within the context of our universe. NASA's Kepler mission has discovered thousands of potential planets around other stars, indicating that Earth is but one of billions of planets in our galaxy. & \textbf{01}\quad 我们正处在人类历史的一个临界点：一个物种正处在获得群星与失去我们称为家园的星球之间。仅仅在过去几年里，我们就极大扩展了关于地球在宇宙背景中位置的知识。NASA 的 Kepler 任务已经在其他恒星周围发现了数千颗潜在行星，这表明地球只是银河系数十亿颗行星中的一颗。 \\ \hline
\textbf{02}\quad Kepler's a space telescope that measures the subtle dimming of stars as planets pass in front of them, blocking just a little bit of that light from reaching us. Kepler's data reveals planet sizes, as well as their distance from their parent star. Together, this helps us understand whether these planets are small and rocky, like the terrestrial planets in our own solar system, and also how much light they receive from their parent sun. In turn, this provides clues as to whether these planets that we discover might be habitable or not. & \textbf{02}\quad Kepler 是一台空间望远镜，它测量恒星在行星从前方经过时出现的微弱变暗，也就是行星挡住了一点点到达我们的星光。Kepler 的数据揭示了行星的大小，以及它们与母恒星的距离。合在一起，这帮助我们理解这些行星是否像我们太阳系中的类地行星一样小而岩质；反过来，这也提供线索，帮助判断我们发现的这些行星是否可能宜居。 \\ \hline
\textbf{03}\quad Unfortunately, at the same time as we're discovering this treasure trove of potentially habitable worlds, our own planet is sagging under the weight of humanity. 2014 was the hottest year on record. Glaciers and sea ice that have been with us for millennia are now disappearing in a matter of decades. These planetary-scale environmental changes that we have set in motion are rapidly outpacing our ability to alter their course. & \textbf{03}\quad 不幸的是，在我们发现这批潜在宜居世界的宝藏的同时，我们自己的星球正在人类的重量下下沉。2014 年是有记录以来最热的一年。陪伴我们数千年的冰川和海冰，如今正在几十年内消失。我们启动的这些行星尺度的环境变化，正在迅速超过我们改变其进程的能力。 \\ \hline
\textbf{04}\quad But I'm not a climate scientist. I'm an astronomer. I study planetary habitability as influenced by stars with the hopes of finding the places in the universe where we might discover life beyond our own planet. You could say that I look for choice alien real estate. Now, as somebody who is deeply embedded in the search for life in the universe, I can tell you that the more you look for planets like Earth, the more you appreciate our own planet itself. Each one of these new worlds invites a comparison between the newly discovered planet and the planets we know best, those of our own solar system. & \textbf{04}\quad 但我不是气候科学家。我是天文学家。我研究受恒星影响的行星宜居性，希望找到宇宙中我们可能发现地球之外生命的地方。你可以说，我是在寻找优质的外星地产。现在，作为一个深深参与宇宙生命搜寻的人，我可以告诉你：你越是寻找像地球一样的行星，就越会珍惜我们自己的星球。每一个新世界都会邀请我们把新发现的行星，与我们最熟悉的行星，也就是我们太阳系中的行星进行比较。 \\ \hline
\textbf{05}\quad Consider our neighbour, Mars. Mars is small and rocky, and though it's a bit far from the sun, it might be considered a potentially habitable world if found by a mission like Kepler. Indeed, it's possible that Mars was habitable in the past. And in part, this is why we study Mars so much. Our rovers, like Curiosity, crawl across its surface, scratching for clues as to the origins of life as we know it. Orbiters like the MAVEN mission sample the Martian atmosphere, trying to understand how Mars might have lost its past habitability. Private space flight companies now offer not just a short trip to near space, but the tantalizing possibility of living our lives on Mars. & \textbf{05}\quad 想想我们的邻居火星。火星小而岩质，虽然离太阳有点远，但如果它是由 Kepler 这样的任务发现的，它可能会被认为是一个潜在宜居世界。事实上，火星过去可能曾经宜居。这也是我们如此大量研究火星的部分原因。我们的探测车，比如 Curiosity，在它的表面爬行，寻找关于我们所知生命起源的线索。像 MAVEN 任务这样的轨道器采样火星大气，试图理解火星可能是怎样失去过去宜居性的。私人航天公司现在提供的不只是近太空短途旅行，还提供了在火星上生活这一诱人的可能性。 \\ \hline
\textbf{06}\quad But though these Martian vistas resemble the deserts of our own home world, places that are tied in our imagination to ideas about pioneering and frontiers, compared to Earth, Mars is a pretty terrible place to live. Consider the extent to which we have not colonized the deserts of our own planet, places that are lush by comparison with Mars. Even in the driest, highest places on Earth, the air is sweet and thick with oxygen, exhaled from thousands of miles away by our rainforests. & \textbf{06}\quad 但是，尽管这些火星景观像我们自己家园世界中的沙漠，而沙漠在我们的想象中又与开拓和边疆这些观念相连，但与地球相比，火星是一个相当糟糕的居住地。想想看，我们甚至还没有殖民自己星球上的沙漠，而这些地方与火星相比已经郁郁葱葱。即使在地球上最干燥、最高的地方，空气也甜美而充满氧气，那是数千英里外我们的雨林呼出的氧气。 \\ \hline
\textbf{07}\quad I worry. I worry that this excitement about colonizing Mars and other planets carries with it a long, dark shadow, the implication and belief by some that Mars will be there to save us from the self-inflicted destruction of the only truly habitable planet we know of, the Earth. As much as I love interplanetary exploration, I deeply disagree with this idea. There are many excellent reasons to go to Mars, but for anyone to tell you that Mars will be there to back up humanity is like the captain of the Titanic telling you that the real party is happening later on the lifeboats. LAUGHTER APPLAUSE Thank you. & \textbf{07}\quad 我担心。我担心这种关于殖民火星和其他行星的兴奋，带着一个漫长而黑暗的影子：它暗示并让一些人相信，火星会在那里拯救我们，让我们逃离对我们所知唯一真正宜居的星球，也就是地球，所造成的自我毁灭。尽管我非常热爱行星际探索，但我深深不同意这种想法。去火星有很多极好的理由，但如果有人告诉你火星会在那里为人类做备份，那就像泰坦尼克号船长告诉你，真正的派对稍后会在救生艇上举行。（笑声）（掌声）谢谢。 \\ \hline
\textbf{08}\quad APPLAUSE But the goals of interplanetary exploration and planetary preservation are not opposed to one another. No, they're in fact two sides of the same goal, to understand, preserve and improve life into the future. The extreme environments of our own world are alien vistas. They're just closer to home. If we can understand how to create and maintain habitable spaces out of hostile, inhospitable spaces here on Earth, perhaps we can meet the needs of both preserving our own environment and moving beyond it. & \textbf{08}\quad （掌声）但是，行星际探索和行星保护的目标并不彼此对立。不，它们实际上是同一个目标的两面：理解、保护并改善走向未来的生命。我们自己世界中的极端环境也是外星景观。它们只是离家更近。如果我们能够理解如何在地球上这些敌对、不宜居的空间中创造并维持宜居空间，也许我们就能同时满足保护自身环境和走出环境之外的需求。 \\ \hline
\textbf{09}\quad I leave you with a final thought experiment, Fermi's Paradox. Many years ago, the physicist Enrico Fermi asked that, given the fact that our universe has been around for a very long time and we expect that there are many planets within it, we should have found evidence for alien life by now. So where are they? Well, one possible solution to Fermi's Paradox is that as civilizations become technologically advanced enough to consider living amongst the stars, they lose sight of how important it is to safeguard the home worlds that fostered that advancement to begin with. & \textbf{09}\quad 最后，我留给你们一个思想实验：费米悖论。许多年前，物理学家 Enrico Fermi 问道，既然我们的宇宙已经存在了很长时间，而且我们预计其中有许多行星，那么到现在我们应该已经发现外星生命的证据了。所以他们在哪里？嗯，费米悖论的一种可能解答是：当文明在技术上发展到足以考虑生活在群星之间时，它们会忽视保护孕育这种进步的母星有多重要。 \\ \hline
\textbf{10}\quad It is hubris to believe that interplanetary colonization alone will save us from ourselves, but planetary preservation and interplanetary exploration can work together. If we truly believe in our ability to bend the hostile environments of Mars for human habitation, then we should be able to surmount the far easier task of preserving the habitability of the Earth. Thank you. & \textbf{10}\quad 认为仅靠行星际殖民就能把我们从自己手中拯救出来，是一种傲慢。但行星保护和行星际探索可以共同进行。如果我们真的相信自己有能力把火星这种敌对环境改造成适合人类居住的地方，那么我们也应该能够克服一个容易得多的任务：保护地球的宜居性。谢谢。 \\ \hline
\end{longtable}


\clearpage
\section{Session 2.5 Activity 4: The Art of Debate (Part 1)}
\addcontentsline{toc}{subsection}{中文译名：辩论的艺术（上）}

\renewcommand{\arraystretch}{1.23}
\begin{longtable}{>{\TranscriptColumn}p{0.47\textwidth} >{\TranscriptColumn}p{0.47\textwidth}}
\rowcolor{headerBg}\textbf{English Transcript} & \textbf{中文翻译} \\ \hline
\textbf{01}\quad When I tell people that I'm on the debate team, I generally get one of two responses. The first being, I could never do that. The second, wow, people have always told me I would be great at debate. I love to argue. Well, both of these responses illustrate a real problem with how we view debate today. Debate is seen as a fight. & \textbf{01}\quad 当我告诉别人我在辩论队时，我通常会得到两种回应。第一种是：“我永远做不到。”第二种是：“哇，别人一直说我会很适合辩论。我很喜欢争论。”好吧，这两种回应都说明了我们今天看待辩论方式中的一个真实问题：辩论被看作一场打斗。 \\ \hline
\textbf{02}\quad If an argument is seen as a fight, then the winner is the person who can talk the loudest, who looks the prettiest on stage, or who can garner the most support. Yet debate is so much more than that. It's about who is the best speaker and who can illustrate their points most strongly. & \textbf{02}\quad 如果一个论证被看作打斗，那么赢家就是那个能说得最大声、在台上看起来最漂亮，或者能争取到最多支持的人。是的，但辩论远不止这些。它关乎谁是最好的演讲者，谁能最有力地说明自己的观点。 \\ \hline
\textbf{03}\quad Now, this isn't all just a shameless plug to try and get you to join the debate team. But you should join. This is a shameless plug to get you to reevaluate the way you form an argument and an opinion, because real debate needs facts, logic, reasoning, and refutation to hold up. & \textbf{03}\quad 现在，这并不全只是我厚着脸皮想让你加入辩论队的宣传。但你确实应该加入。这是一段厚着脸皮的宣传，目的是让你重新评估自己形成论证和观点的方式，因为真正的辩论需要事实、逻辑、推理和反驳来支撑。 \\ \hline
\textbf{04}\quad So, I'm going to tell you the story of the first time I remember losing a debate. So, I was with my mom and my younger sister in the car, or my older sister in the car, and we were driving home from school. And I told them what I learned that day, that Mount Everest is the tallest mountain in the whole world. My sister immediately shot me down. She said, actually, Mauna Kea in Hawaii is way taller than that. I started screaming at her. I said, no, you're wrong. A little spoiler alert, I was wrong. & \textbf{04}\quad 所以，我要告诉你我记得自己第一次输掉辩论的故事。当时我和妈妈、妹妹，或者说姐姐，坐在车里，我们正从学校回家。我告诉她们我那天学到的东西：珠穆朗玛峰是全世界最高的山。我姐姐立刻驳倒了我。她说，其实夏威夷的 Mauna Kea 比它高得多。我开始冲她大喊。我说：“不，你错了。”小小剧透一下：错的是我。 \\ \hline
\textbf{05}\quad So, actually, Mauna Kea is a mountain in Hawaii, but it's like a volcano. And if you go below sea level, it's about a mile taller than Everest. I didn't know that at the time. And I had learned from my teacher that Mount Everest was the tallest, so I believed it. And I fought to defend that point, just with words. My sister would win, but in a fight. But I fought to defend that point anyway, even though I didn't have all the information. & \textbf{05}\quad 所以，实际上，Mauna Kea 是夏威夷的一座山，但它像一座火山。如果从海平面以下算起，它大约比珠穆朗玛峰高一英里。那时我并不知道这一点。我从老师那里学到珠穆朗玛峰最高，所以我相信这一点。我用语言努力维护这个观点。我的姐姐会赢，但那是一场打斗。而我是在没有掌握所有信息的情况下，努力把自己的观点拼起来。 \\ \hline
\end{longtable}


\clearpage
\section{Session 2.5 Activity 5: The Art of Debate (Part 2)}
\addcontentsline{toc}{subsection}{中文译名：辩论的艺术（下）}

\renewcommand{\arraystretch}{1.23}
\begin{longtable}{>{\TranscriptColumn}p{0.47\textwidth} >{\TranscriptColumn}p{0.47\textwidth}}
\rowcolor{headerBg}\textbf{English Transcript} & \textbf{中文翻译} \\ \hline
\textbf{01}\quad So how can you avoid looking like a complete idiot in an argument? Well, you need to learn how to build an informed argument. And the first step to that is understanding the question and the context to the question. So if we look back to the mountain example, if we had posed the question, what is the tallest mountain including below and above sea level heights, I would have had the chance to think, to examine my position. Because I would have thought that, does my statistic apply to above or just below sea level? And I could have understood that in the context, my statistic was irrelevant. And there were other possibilities to what the answer of the question could have been. & \textbf{01}\quad 那么，你怎样才能避免在争论中看起来像个彻头彻尾的傻瓜？你需要学会建立一个有信息支持的论证。第一步是理解问题，以及这个问题的语境。所以，如果我们回到山的例子，如果我们提出的问题是“包括海平面以下和海平面以上的高度在内，最高的山是什么”，我就会有机会思考，审视自己的立场。因为我会去想：我的统计数据适用于海平面以上，还是也适用于海平面以下？我也会理解，在这个语境中，我的数据并不相关；这个问题的答案还可能有其他可能性。 \\ \hline
\textbf{02}\quad And this doesn't just apply to mountains, because when you are building a point in a discussion, you need to know what you're saying. You need to understand what the question being asked is. For this question, I was led astray; I didn't understand if we were including below and above sea level, so I wasn't able to build an informed argument. And as we come to topics that people are more passionate about, people more often forget to look at the reasoning and the context behind the argument, which is very important to formulating and defending an opinion. & \textbf{02}\quad 这并不只适用于山。因为当你在讨论中建立一个观点时，你需要知道自己在说什么。你需要理解被问到的问题是什么。对于这个问题，我没有理解我们是否把海平面以下和以上都包括在内，所以我无法建立一个有信息支持的论证。而当我们讨论人们更有热情的话题时，人们更常忘记去看论证背后的推理和语境，而这对于形成并捍卫一个观点非常重要。 \\ \hline
\textbf{03}\quad Secondly, you need to understand the argument from multiple and all relevant perspectives. So, if we go back to the Mauna Kea example again, I, Mount Everest was the only mountain that I knew the height of. Not the height, but I knew that it was the tallest. So it was the only mountain that I had heard of. But, if I had examined other mountains, I probably would have come across mountains whose height was included below and above sea level. And in that case, I would have been able to reevaluate my argument in the context of the question and understand why my sister thought that I was wrong. & \textbf{03}\quad 第二，你需要从多个相关视角，或者说所有相关视角，来理解这个论证。所以，如果我们再次回到 Mauna Kea 的例子，我只知道珠穆朗玛峰这一座山的高度。不是高度，准确地说，是我知道它是最高的。所以它是我听说过的唯一一座山。但是，如果我考察过其他山，我可能就会遇到一些山，它们的高度包括海平面以下和海平面以上。在那种情况下，我就能够在问题的语境中重新评估自己的论证，并理解为什么我姐姐认为我错了。 \\ \hline
\textbf{04}\quad Furthermore, this applies to lots of arguments. Often people come across as ignorant in a discussion because they have not considered all the people who are involved in the argument. There are lots of opinions different from yours, and even an argument that can seem black and white in your opinion can be seen as very colorful in someone else's. So, in order to form an argument, you need to look at all the relevant perspectives. & \textbf{04}\quad 此外，这也适用于很多论证。人们在讨论中常常显得无知，是因为他们没有考虑所有卷入这个论证的人。每个人，世界上有很多与你不同的意见；即使一个论证在你看来非黑即白，在别人那里也可能被看作色彩丰富。所以，为了形成一个论证，你需要看所有相关的视角。 \\ \hline
\textbf{05}\quad And lastly, and most importantly, in order to form an argument, you need to look at the side that's opposing yours. This is extremely important. In debate, we call this the con to your pro, or the affirmative to your negative. This person inherently thinks you're wrong. The fact that there is a debate means that there's a person who thinks that you're wrong. And this is very important because this person has a reason why they think they're right. They think this reason is very right and they are going to defend it. & \textbf{05}\quad 最后，也是最重要的一点，为了形成一个论证，你需要看与你相反的那一方。这极其重要。在辩论中，我们把它叫作你的正方对应的反方，或者你的反方对应的正方。这个人从根本上认为你错了。有辩论这一事实，就意味着有人认为你错了。而这非常重要，因为这个人有理由认为自己是对的。他们认为这个理由非常正确，并且会为它辩护。 \\ \hline
\textbf{06}\quad Yet, if you look at the reason why you're right and discover why they are wrong, you will be able to prove your point in an argument. So, if we go back to the mountain example again, we can see that my sister thought she was right because she was looking at the statistic from below and above sea level. Yet, if I rephrased my argument to be that Mount Everest is the tallest mountain in the world above sea level, I would win that argument every time. & \textbf{06}\quad 然而，如果你看清自己为什么正确，也发现他们为什么错误，你就能够在争论中证明自己的观点。所以，如果我们再次回到山的例子，我们可以看到，我姐姐认为自己正确，是因为她是从海平面以下和海平面以上的统计数据来看问题的。然而，如果我把自己的论证改写为“珠穆朗玛峰是世界上海平面以上最高的山”，我每一次都会赢得这个论证。 \\ \hline
\textbf{07}\quad In the status quo, when we debate, we make quick judgments and avidly defend them, even though we later often come to find that we were misinformed and answered wrongly. Real debate requires you to slow down and become a better consumer of information in order to prove a point in an argument. In order to do this, you need to look at all perspectives and contexts in order to answer the question. You also need to look at the question in context, in addition to looking at the opposing side of the argument. If you do all these, you will learn to debate what's in your heart with your mind instead of letting your heart dictate how your mind perceives information. & \textbf{07}\quad 在现状中，当我们辩论时，我们会使用快速判断，并热切地捍卫这些判断，尽管后来我们常常发现自己被误导了，回答错了。真正的辩论要求你慢下来，成为一个更好的信息消费者，这样你才能在争论中证明一个观点。为了做到这一点，你需要看所有视角和语境，才能回答问题。除了看反方观点，你还需要把问题放在语境中来看。如果你做到这些，你就会学会用头脑去辩论心中所想，而不是让内心情绪决定你的头脑如何理解信息。 \\ \hline
\textbf{08}\quad So, the next time you're in an argument or a discussion and you're getting frazzled, I implore you, don't raise your voice but improve your argument. Thank you so much. Join the debate team. & \textbf{08}\quad 所以，下一次当你陷入争论或讨论，并且开始慌乱时，我恳请你：不要提高嗓门，而要改进你的论证。非常感谢。加入辩论队吧。 \\ \hline
\end{longtable}


\clearpage
\section{Session 2.6 Activity 7 \& 8: How Our Food Choices Affect Climate Change}
\addcontentsline{toc}{subsection}{中文译名：我们的食物选择如何影响气候变化}

\renewcommand{\arraystretch}{1.23}
\begin{longtable}{>{\TranscriptColumn}p{0.47\textwidth} >{\TranscriptColumn}p{0.47\textwidth}}
\rowcolor{headerBg}\textbf{English Transcript} & \textbf{中文翻译} \\ \hline
\textbf{01}\quad I have to admit that when it comes to food, I'm a total sucker. Whether it's sugar or grease or carbs, pretty much bring it on. And I spend a lot of time in Montana. So for me, that medium-rare grass-fed rib-eye steak? Pretty much as good as it gets. I know, I don't do it often. And when I do, I gotta admit, I feel a little conflicted. And that's for a lot of reasons, including the planet. & \textbf{01}\quad 我必须承认，说到食物，我完全经不起诱惑。不管是糖、油脂还是碳水，基本上都来吧。而且我有很多时间待在蒙大拿州。所以对我来说，那块五分熟的草饲肋眼牛排？几乎就是美味的顶点。我知道，我并不经常吃。但当我吃的时候，我得承认，我会有一点矛盾。这有很多原因，其中也包括地球。 \\ \hline
\textbf{02}\quad But how big of a problem is what I eat? I mean, does it really make much of a dent in something as huge as global warming? But it turns out what we put on our plates matters a lot. About 25\% of all the global climate change problems we're seeing can be attributed back to the food and the choices that we're actually making about what we eat on a daily basis. This is greater than all the cars on the planet. In fact, it's about twice as much global warming pollution as the cars. & \textbf{02}\quad 但我吃什么，到底是多大的问题？我是说，它真的会对全球变暖这样巨大的事情产生多少影响吗？事实证明，我们盘子里的东西非常重要。我们正在看到的所有全球气候变化问题中，大约 25\% 可以追溯到食物，以及我们每天实际做出的关于吃什么的选择。这比地球上所有汽车造成的影响还要大。事实上，食物造成的全球变暖污染大约是汽车的两倍。 \\ \hline
\textbf{03}\quad Ben Houlton and Maya Almaraz study the connection between climate and diet at the University of California, Davis. They track how the way we produce food creates greenhouse gases that contribute to global warming. With their data, the team has crunched the numbers to figure out how much carbon pollution is produced by different foods and different diets. A lot of people feel really helpless when it comes to climate change, like they can't make a difference. And what our research is showing is that your personal decisions really can have a big impact. & \textbf{03}\quad Ben Houlton 和 Maya Almaraz 在加州大学戴维斯分校研究气候与饮食之间的联系。他们追踪我们生产食物的方式如何产生温室气体，并因此推动全球变暖。利用他们的数据，这个团队计算了不同食物和不同饮食会产生多少碳污染。很多人在面对气候变化时感到非常无力，觉得自己无法产生影响。而我们的研究显示，你的个人决定确实可以产生巨大影响。 \\ \hline
\textbf{04}\quad So take that grass-fed ribeye steak I love. If you really look at everything that went into making a single serving of beef, you end up emitting around 330 grams of carbon. That's like driving a car three miles. Now, if I choose to have chicken instead, there's more than a five-fold drop in emissions. Switch to fish and you see the number go down even more. Now, look at veggies. If I swapped beef out entirely for lentils, well, I'm down to practically nothing. & \textbf{04}\quad 就拿我喜欢的那块草饲肋眼牛排来说。如果你真的看清做出一份牛肉背后所有投入的东西，最终会排放大约 330 克碳。这就像开车 3 英里。现在，如果我选择吃鸡肉，排放量会下降五倍以上。换成鱼，你会看到数字降得更多。再看看蔬菜。如果我完全用扁豆替代牛肉，那么排放几乎降到没有。 \\ \hline
\textbf{05}\quad So why does beef, and lamb too for that matter, pack such a powerful punch to the planet? Livestock accounts for a little over 14\% of global greenhouse gas emissions. If that sort of seems low to you, consider it's about equal to transportation. We're talking all the cars, trucks, planes, trains, and ships on the planet combined. This is partly because ruminant animals like cows and sheep, they're just gassy. And the methane they produce is at least 25 times more potent than carbon dioxide. Plus, it takes a lot of land, fertilizer, and about a billion tons of grain to feed all that livestock. And you could feed 3.5 billion people with that grain. If we were just directly eating these grains ourselves, it would eliminate a lot of the CO2 that's emitted from cattle production. & \textbf{05}\quad 那么，为什么牛肉，以及羊肉，在这件事上对地球有这么强的冲击？畜牧业占全球温室气体排放的 14\% 多一点。如果这个数字在你看来似乎不高，那就想想，它大约等于交通运输业。我们说的是地球上所有汽车、卡车、飞机、火车和船舶加在一起。这部分原因在于，像牛和羊这样的反刍动物就是会产生很多气体。它们产生的甲烷至少比二氧化碳强 25 倍。此外，饲养所有牲畜需要大量土地、化肥和大约 10 亿吨谷物。而这些谷物可以养活 35 亿人。如果我们自己直接吃这些谷物，就会消除很多来自牛肉生产的二氧化碳排放。 \\ \hline
\textbf{06}\quad So it's clear that meat has a pretty big carbon load, but it's also worth remembering that not all livestock is raised equally. In parts of the American West, for instance, ranchers are working to raise livestock in ways that actually help restore the land. And they're experimenting with ways that soil and grasslands can be used to keep carbon pollution out of the air. But even these sustainable ranchers will really tell you, we're probably eating too much meat. I know a lot of people who, if you don't serve them meat for lunch or dinner, they're kind of like, well, when is the meat coming out? It's to the point now where the US actually has one of the highest meat footprints per capita. & \textbf{06}\quad 所以很清楚，肉类有相当大的碳负荷。但也值得记住，并非所有牲畜都是以同样方式饲养的。例如，在美国西部的一些地区，牧场主正在努力以实际上有助于恢复土地的方式饲养牲畜。他们也在尝试利用土壤和草地把碳污染留在空气之外的方法。但即使是这些可持续牧场主也会真的告诉你：我们可能吃了太多肉。我认识很多人，如果午餐或晚餐不给他们上肉，他们就会有点像：“肉什么时候来？” 现在已经到了这样的程度：美国实际上拥有世界上最高的人均肉类足迹之一。 \\ \hline
\textbf{07}\quad So what about not eating meat at all? Vegan is the way to go for the least impact on the planet. But it's not that much different in terms of emissions than, say, a vegetarian diet. And the team found that the environmental impact of the Mediterranean diet is pretty similar to vegan and vegetarian diets. It's a lot less meat heavy than what Americans are used to. So fish and poultry a few times a week, beef maybe once a month, plenty of plant-based foods, and of course, loads of olive oil. & \textbf{07}\quad 那么完全不吃肉怎么样？从对地球影响最小的角度看，纯素饮食是方向。但在排放方面，它和比如素食饮食并没有差很多。团队还发现，地中海饮食的环境影响与纯素和素食饮食相当接近。它比美国人习惯的饮食少得多肉。所以，每周吃几次鱼和禽肉，牛肉也许每月吃一次，大量植物性食物，当然还有很多橄榄油。 \\ \hline
\textbf{08}\quad Eliminating, like, 90\% of your meat intake is more important than eliminating all of your meat. We don't all have to be vegan. We don't all even have to be vegetarian. If we can just reduce our meat intake, every little bit helps. And if you can bring it down a lot, you can help the climate a lot. If we all just switched to a Mediterranean diet, it could actually solve 15\% of global warming pollution by 2050. If everyone were to move toward it, that is equivalent to taking somewhere around a billion cars off of the streets in terms of vehicle emissions each year. So that kind of a footprint is big time. & \textbf{08}\quad 减少大约 90\% 的肉类摄入，比消除全部肉类摄入更重要。我们不必人人都成为纯素者。我们甚至不必人人都成为素食者。如果我们只是减少肉类摄入，每一点都有帮助。如果你能大幅减少，就能对气候有很大帮助。如果我们都只是转向地中海饮食，到 2050 年，它实际上可以解决 15\% 的全球变暖污染。如果每个人都朝这个方向移动，就车辆排放而言，这相当于让大约 10 亿辆汽车离开街道。所以这种足迹是非常大的。 \\ \hline
\textbf{09}\quad But say you still want more meat than the Mediterranean diet recommends. Just cutting down your portion size to the doctor-recommended four ounces can reduce your emissions by half. That's huge. In fact, the doctors are telling us we're eating about twice as much meat as we really need for a healthy diet. The good news is we are listening to our doctors. In the last decade, there's been a 19\% drop in the amount of beef we eat. All these things that you're already being told are good for you also happen to be good for the planet. & \textbf{09}\quad 但假设你仍然想吃比地中海饮食建议更多的肉。仅仅把份量减少到医生建议的 4 盎司，就能把你的排放减少一半。这非常大。事实上，医生告诉我们，为了健康饮食，我们吃的肉大约是实际需要的两倍。好消息是，我们正在听医生的话。在过去十年里，我们吃的牛肉数量下降了 19\%。所有这些你已经被告知对自己有益的事情，碰巧也对地球有益。 \\ \hline
\textbf{10}\quad So what we eat really is a big part of the climate puzzle. I mean, we may not all be able to afford an electric car or putting solar panels on our house, but we all have to eat every day. And these choices we make can add up to really big numbers. And since meat has a pretty big carbon load, we need to be thoughtful about how much we eat. As for that ribeye steak that I really love, I am honestly trying to cut back. Maybe it's just a smaller piece of steak or simply swapping out a meat dish for a veggie burger. It may seem like a small thing, but it really does add up to big impacts. & \textbf{10}\quad 所以，我们吃什么确实是气候拼图中的一大部分。我的意思是，我们也许并不是所有人都买得起电动车，或者在房子上安装太阳能板，但我们每个人每天都必须吃饭。而我们做出的这些选择可以累积成非常大的数字。既然肉类有相当大的碳负荷，我们就需要认真思考自己吃多少肉。至于那块我真的很爱的肋眼牛排，我真诚地在努力减少。也许只是吃一小块牛排，或者只是把一道肉菜换成素食汉堡。这看起来像一件小事，但确实会累积成很大的影响。 \\ \hline
\end{longtable}


\end{document}
